\documentclass[10pt,a4paper]{amsart}
\usepackage[margin=19mm]{geometry}
\usepackage{amsmath,amssymb,mathtools}
\usepackage[scr=rsfs,cal=euler]{mathalfa}
\usepackage{xfrac}
\usepackage[unicode,hidelinks,bookmarks=false]{hyperref}
\newcommand{\ie}{i.e.\ }
\newcommand{\cf}{cf.\ }
\newcommand{\resp}{respectively\ }
\newcommand{\Subsection}{\subsection}
\providecommand{\nobreakdash}{\nobreak\hbox{-}}
\setlength{\emergencystretch}{2em}
\multlinegap0pt
\usepackage{xcolor}
\usepackage{tikz}
\usepackage{amssymb}
\newtheorem{theorem}{Theorem}
\newtheorem{corollary}{Corollary}
\newcommand{\N}{\mathbb N}
\newcommand{\eps}{\varepsilon}
\newcommand{\set}[1]{\left\{#1\right\}}
\title{On the escape rate for intermittent maps with holes shrinking around the indifferent fixed point}
\author{Claudio Bonanno, Sharvari Neetin Tikekar}
\date{Source: arXiv:2601.15908v3 (CC BY 4.0)}
\begin{document}
\maketitle

\noindent\emph{Source and licence.} On the escape rate for intermittent maps with holes shrinking around the indifferent fixed point. Version arXiv:2601.15908v3 (CC BY 4.0), distributed by arXiv under the Creative Commons Attribution 4.0 International licence, http://creativecommons.org/licenses/by/4.0/, as recorded in the arXiv licence metadata for this exact version and shown on the abstract page https://arxiv.org/abs/2601.15908v3. Full licence notice: \texttt{licenses/arxiv-2601.15908-CC-BY-4.0.txt}.\par\smallskip

\noindent\textit{Editorial scope.} Counted unit: the article's corollary general-hole (source0.tex lines 668--674). The article's complete original proof of it is delivered with it (source0.tex lines 676--705, one \texttt{proof} environment of 30 lines). No mathematics is written, repaired or re-derived: the statement and the proof are the article's own lines, in the article's own notation, and no retained line is substituted or re-flowed.  Retained uncounted as the source-local context the proof needs: lines 168--170; lines 376--378; lines 638--644; lines 640--642; lines 650--652.  Nothing was omitted from any retained span, so the package carries no omission record.  No retained line contains a citation, so the wrapper carries no bibliography and no bibliography entry.  No retained line contains a figure environment, an image-inclusion command or drawing code, so no image is delivered and the package declares no asset dependency.  Editorial formatting only, in addition: an AMS \texttt{amsart} preamble that restates the article's own declarations and macros used by the retained lines, namely the environments \texttt{theorem}, \texttt{corollary} on separate counters, together with the standard packages those lines need.  The retained environments print in this document as Theorem 1, Corollary 1 in order of appearance, whereas the article prints the same items with the numbering of its own full document.  The complete native source of the article as distributed in the arXiv e-print for this exact version is bundled unchanged as \texttt{sources/193198/source0.tex}.
\begin{equation} \label{preimages}
        a_{0}:=1, \ \ \ a_{1}:= a, \ \ \ \text{and} \ \ \   a_{n}:=\phi_{0}(a_{n-1}), \  n \ge 2. 
\end{equation} 

\begin{equation}\label{escape-rate-F}
\gamma_\mu(H) := \limsup_{n\to \infty}\, \frac {-\log \Big( \mu\set{x\in [0,1]\, :\, F^j(x) \not\in H\, ,\, \forall\, j=0,\dots, n-1} \Big)}{n}\, .
\end{equation}

\begin{theorem}\label{thm:shrinking}
Consider the parabolic dynamical system on the interval satisfying (H1)-(H5) and (W1)-(W2), and fix a Markov hole $H:=[0,a_N]$ for $N \ge 2$, with $a_N$ defined as in \eqref{preimages}. Then, as $N\to \infty$, or equivalently as $m(H)\to 0^+$ with $m$ the Lebesgue measure, we have
\begin{equation}\label{final-formula}
\gamma_\mu(H) \sim \left\{ \begin{array}{ll} c_s \cdot m(H)\, , & \text{if $s\in (0,1)$;}\\[0.2cm] c_s \cdot \frac{m(H)}{-\log m(H)}\, , & \text{if $s=1$;}\\[0.2cm] c_s \cdot m(H)^s\, , & \text{if $s>1$,} \end{array} \right.
\end{equation}
where $c_s>0$ denotes a constant only depending on the system and not on the hole.
\end{theorem}

\begin{equation}
\gamma_\mu(H) \sim \left\{ \begin{array}{ll} c_s \cdot m(H)\, , & \text{if $s\in (0,1)$;}\\[0.2cm] c_s \cdot \frac{m(H)}{-\log m(H)}\, , & \text{if $s=1$;}\\[0.2cm] c_s \cdot m(H)^s\, , & \text{if $s>1$,} \end{array} \right.
\end{equation}

\begin{equation}\label{asymp-an}
m(H) = a_N \sim const \cdot N^{-\frac 1s}\, ,\qquad \rho(A_N) \sim const \cdot N^{-1-\frac 1s}
\end{equation}

\begin{corollary} \label{general-hole}
Consider the parabolic dynamical system on the interval satisfying (H1)-(H5) and (W1)-(W2), and fix a hole $H_\eps=[0,\eps]$ with $\eps>0$. Then, as $\eps\to 0^+$,
\[
\gamma_\mu(H_\eps) \sim \left\{ \begin{array}{ll} c_s \cdot m(H_\eps)\, , & \text{if $s\in (0,1)$;}\\[0.2cm] c_s \cdot \frac{m(H_\eps)}{-\log m(H_\eps)}\, , & \text{if $s=1$;}\\[0.2cm] c_s \cdot m(H_\eps)^s\, , & \text{if $s>1$,} \end{array} \right.
\]
where $c_s$ is a constant as in Theorem \ref{thm:shrinking}.
\end{corollary}

\begin{proof}
The proof relies on the observation that the asymptotic behaviour of escape rate for general holes $H_\eps$ is controlled by the asymptotic behaviour of escape rate of Markov holes as we describe below. It is immediate that for all $\eps>0$ there exists $N_\eps \in \N$ such that $a_{N_\eps+1} < \eps \le a_{N_\eps}$. Hence, we find two Markov holes, $H_{N_\eps+1} = [0,a_{N_\eps+1}]$ and $H_{N_\eps} = [0,a_{N_\eps}]$, such that $H_{N_\eps+1} \subset H_\eps \subseteq H_{N_\eps}$. 

The definition of the escape rate \eqref{escape-rate-F} implies that $\gamma_\mu(H_{N_\eps+1})\le \gamma_\mu(H_\eps)\le \gamma_\mu(H_{N_\eps})$. Hence, we can use \eqref{final-formula} to obtain a result about the behaviour of $\gamma_\mu(H_\eps)$ as $\eps\to 0^+$.

By \eqref{asymp-an}, as $N\to \infty$,
\[
\frac{m(H_{N+1})}{m(H_N)} = \frac{a_{N+1}}{a_N} \sim \Big( \frac{N}{N+1} \Big)^{\frac 1s} \sim 1\, .
\]

Let's consider first the case $s> 1$. For any $\eps >0$, we have,
\[
 \frac{\gamma_\mu(H_{N_\eps+1})}{c_s\cdot m(H_{N_\eps+1})^s}\left[ \frac{m(H_{N_\eps+1})}{m(H_{N_\eps})} \right]^s \le \frac{\gamma_\mu(H_\eps)}{c_s \cdot m(H_\eps)^s} \le \frac{\gamma_\mu(H_{N_\eps})}{c_s\cdot m(H_{N_\eps})^s} \left[ \frac{m(H_{N_\eps})}{m(H_{N_\eps +1})} \right]^s .
\]
%\le C\, \Big( \frac{a_{N_\eps}}{\eps} \Big)^s \le C\, \Big( \frac{a_{N_\eps}}{a_{N_\eps+1}} \Big)^s\, 
%\frac 1C\, \Big( \frac{a_{N_\eps+1}}{a_{N_\eps}} \Big)^s \le \frac 1C\, \Big( \frac{a_{N_\eps+1}}{\eps} \Big)^s \le

Then for $\eps$ small enough, hence $N_\eps$ big enough, the lower and upper bounds above are $\sim 1$. Hence using the sandwich theorem we get the desired asymptotic.
 

For $s=1$, the same argument leads to 
\[
\frac{\gamma_\mu(H_\eps)}{c_s \cdot m(H_\eps)/\log \frac{1}{m(H_\eps)}} \le \frac{\gamma_\mu(H_{N_\eps})}{c_s \cdot \left( m(H_{N_\eps}) / \log \frac{1}{m(H_{N_\eps})} \right)} \cdot \frac{\left[ m(H_{N_\eps}) / \log \frac{1}{m(H_{N_\eps})}\right]}{\left[ m(H_{N_\eps+1}) / \log \frac{1}{m(H_{N_\eps+1})}\right]} \sim 1 \, ,
\]
and,
\[
\frac{\gamma_\mu(H_\eps)}{c_s \cdot m(H_\eps)/\log \frac{1}{m(H_\eps)}} \ge \frac{\gamma_\mu(H_{N_\eps+1})}{c_s \cdot \left( m(H_{N_\eps+1}) / \log \frac{1}{m(H_{N_\eps+1})} \right)} \cdot \frac{\left[ m(H_{N_\eps+1}) / \log \frac{1}{m(H_{N_\eps+1})}\right]}{\left[ m(H_{N_\eps}) / \log \frac{1}{m(H_{N_\eps})}\right]} \sim 1 \,.
\]
The argument for $s\in (0,1)$ follows along the same lines, hence, concludes the proof.
\end{proof}

\end{document}
